% ##############################################################################
\section{Selection, scheduling, and probing}
\label{sec:probing}

This section describes how an active contract receives its probes: which
canaries, at what times, through which path, and with which failure semantics.

% ==============================================================================
\subsection{Keyed, family-balanced selection}
\label{sec:selection}

Let $\bank^{*}$ be the set of \texttt{ACTIVE} canaries, grouped by family into
$\bank^{*}_f$, and let $N$ be the number of probes of contract $\kappa$ (the
contract field \texttt{canaries\_target}, default 10). Let $k$ be the platform
secret. The selection procedure derives a seed from the contract identifier,
\begin{equation}
  \label{eq:selection_seed}
  s(\kappa) = \operatorname{int}\!\left(
    \operatorname{HMAC\text{-}SHA256}_k(\text{decimal}(\mathit{id}_\kappa))[0{:}8]
  \right),
\end{equation}
that is, the first 64 bits of the keyed hash \cite{rfc2104}. It seeds a
pseudo-random generator, shuffles the canaries of each family and the order of
the families, and then draws one canary from each family in turn until $N$
canaries are selected, skipping families that have run out.

\begin{proposition}
\label{prop:balance}
If every family has at least $\lceil N/|\family| \rceil$ canaries in
$\bank^{*}$, the selected batch contains between $\lfloor N/|\family| \rfloor$
and $\lceil N/|\family| \rceil$ canaries of each family, and it contains no
canary twice. For fixed $k$ and $\mathit{id}_\kappa$ the batch is identical on
every invocation.
\end{proposition}

\begin{proof}
Each round of the loop takes one canary from each non-empty family, in a fixed
order, so after $r$ complete rounds each family has contributed $r$ canaries.
The loop stops after $N$ draws, within a round that has taken at most one
canary from each family. Canaries are popped from per-family lists, so none is
drawn twice. The generator is seeded deterministically.
\end{proof}

With six families and $N = 10$, a contract receives one or two canaries per
family. Balance limits the dependence of a contract's evidence on the
difficulty of any single task family. The selection also satisfies property P4
for a party without $k$, subject to the caveat in
Section~\ref{sec:limitations} that the generator after the keyed seed is not a
cryptographic one.

% ==============================================================================
\subsection{Keyed probe schedule}
\label{sec:schedule}

Let the contract window start at $t_0$ with maximum length $W$ seconds. The
schedule divides the window into $N$ equal buckets of width $b = W/N$ and
places probe $i \in \{1,\dots,N\}$ inside bucket $i$. The position within the
bucket is a keyed value,
\begin{equation}
  \label{eq:schedule}
  u_i = \frac{\operatorname{int}\!\left(
      \operatorname{HMAC\text{-}SHA256}_k(\texttt{probe-schedule:}\mathit{id}_\kappa\texttt{:}i)[0{:}8]
    \right)}{2^{64}},
  \qquad
  t_i = t_0 + (i - 1 + 0.2 + 0.6\,u_i)\, b,
\end{equation}
so that $t_i$ lies in the middle $60\%$ of its bucket. Figure~\ref{fig:schedule}
illustrates the construction with five probes.

% rendered_images:begin
% ```raw_latex
% \documentclass[tikz]{standalone}
% \usepackage{amsmath}
% \usetikzlibrary{arrows.meta, decorations.pathreplacing, positioning, calc}
% \definecolor{BandBlue}{RGB}{166,200,244}
% \definecolor{ProbeRed}{RGB}{214,69,69}
% \definecolor{GridGray}{RGB}{140,140,140}
% \begin{document}
% \begin{tikzpicture}[
%   x=2.3cm, y=1cm, font=\sffamily\small,
%   bucket/.style={draw=GridGray, line width=0.4pt, fill=gray!6},
%   band/.style={fill=BandBlue!45, draw=none},
%   probe/.style={circle, fill=ProbeRed, draw=none, minimum size=5pt, inner sep=0pt},
%   stem/.style={draw=ProbeRed, line width=0.8pt},
%   lab/.style={font=\sffamily\footnotesize, align=center}
% ]
%   % Buckets and their eligible middle bands.
%   \foreach \i in {0,1,2,3,4} {
%     \draw[bucket] (\i,0) rectangle (\i+1,1);
%     \fill[band] (\i+0.2,0) rectangle (\i+0.8,1);
%     \draw[GridGray, line width=0.4pt] (\i+0.2,0) -- (\i+0.2,1);
%     \draw[GridGray, line width=0.4pt] (\i+0.8,0) -- (\i+0.8,1);
%   }
%   % Probes at secret positions inside each band: offset (i - 1 + 0.2 + 0.6 u_i) b.
%   \foreach \i/\p/\fam in {0/0.31/{table}, 1/0.72/{arithmetic}, 2/0.45/{string}, 3/0.26/{modular}, 4/0.64/{ordering}} {
%     \node[probe] (p\i) at (\i+\p,0.5) {};
%     \draw[stem] (\i+\p,0.5) -- (\i+\p,1.55);
%     \node[lab, above] at (\i+\p,1.55) {probe \the\numexpr\i+1\relax\\[-1pt]\textit{\fam}};
%   }
%   % Time axis.
%   \draw[-{Latex[length=2.2mm]}, line width=0.6pt] (0,-0.35) -- (5.15,-0.35);
%   \foreach \i/\l in {0/{$t_0$}, 1/{$t_0{+}b$}, 2/{$t_0{+}2b$}, 3/{$t_0{+}3b$}, 4/{$t_0{+}4b$}, 5/{$t_0{+}W$}} {
%     \draw[line width=0.5pt] (\i,-0.28) -- (\i,-0.42);
%     \node[lab, below] at (\i,-0.42) {\l};
%   }
%   % Bucket width brace.
%   \draw[decorate, decoration={brace, mirror, amplitude=4pt}, line width=0.5pt]
%     (2,-1.05) -- (3,-1.05) node[midway, below=5pt, lab] {bucket width $b = W/N$};
%   % Legend.
%   \node[draw=GridGray, line width=0.4pt, fill=BandBlue!45, minimum size=3mm, inner sep=0pt] (l1) at (0.05,-2.05) {};
%   \node[lab, right=2pt of l1, anchor=west] {eligible band: middle $60\%$ of a bucket};
%   \node[draw=GridGray, line width=0.4pt, fill=gray!6, minimum size=3mm, inner sep=0pt] (l2) at (2.45,-2.05) {};
%   \node[lab, right=2pt of l2, anchor=west] {unused margin: $20\%$ at each end};
%   \node[probe] (l3) at (4.55,-2.05) {};
%   \node[lab, right=2pt of l3, anchor=west] {secret probe time};
% \end{tikzpicture}
% \end{document}
% ```
% label=fig:schedule
% caption=Probe times for $N = 5$. Each probe falls at a keyed position in the middle 60\% of its bucket. The family labels are illustrative.
% rendered_images:end
% render_images:begin
\begin{figure}[H]
  \includegraphics[width=\linewidth]{06_probing.tex.figs/06_probing.1.png}
  \caption{Probe times for $N = 5$. Each probe falls at a keyed position in the middle 60\% of its bucket. The family labels are illustrative.}
  \label{fig:schedule}
\end{figure}
% render_images:end

\begin{proposition}
\label{prop:spacing}
All probe times satisfy $t_0 + 0.2\,b \leq t_i \leq t_0 + W - 0.2\,b$, and
consecutive probes satisfy $0.4\,b \leq t_{i+1} - t_i \leq 1.6\,b$.
\end{proposition}

\begin{proof}
From \eqref{eq:schedule}, $t_i = t_0 + (i - 1 + \beta_i)\,b$ with $\beta_i = 0.2 +
0.6\,u_i \in [0.2, 0.8]$. The first bound follows from $i = 1$ and $i = N$.
Subtracting consecutive times gives $t_{i+1} - t_i = (1 + \beta_{i+1} -
\beta_i)\,b$, and $\beta_{i+1} - \beta_i \in [-0.6, 0.6]$.
\end{proof}

The margins keep probes away from the window boundaries, and the spacing bound
prevents both long gaps and bursts. For the demo window of 60 s and
$N = 10$, $b = 6$~s and consecutive probes are between 2.4~s and 9.6~s apart.
For a window of 300 s, $b = 30$~s. The schedule is a pure function of the
contract, the active canaries, and $k$, so it can be recomputed and audited. The
function \texttt{get\_due\_probes} returns the slots with $t_i \leq \text{now}$
that are not yet claimed, in chronological order.

% ==============================================================================
\subsection{Executing a probe}
\label{sec:execution}

One scheduler pass, \texttt{run\_due\_probes}, loads the active contracts and
the active canaries, rebuilds each contract's schedule, subtracts the claimed
slots, and executes each due slot. Figure~\ref{fig:probe-execution} shows the
steps of one probe.

% rendered_images:begin
% ```graphviz
% digraph ProbeExecution {
%   // state change : teal | process : blue | external : violet | failure : coral
%   bgcolor="transparent";
%   pad="0.10";
%   splines=spline;
%   nodesep=0.25;
%   ranksep=0.38;
%   node [shape=box, style="rounded,filled", fillcolor="#FFFFFF", color="#9AA9B8", penwidth=1.8, fontname="Helvetica", fontcolor="#243B53", fontsize=11, margin="0.14,0.08", height=0.42];
%   edge [color="#A3B1C0", penwidth=1.3, arrowhead=vee, arrowsize=0.75, fontname="Helvetica", fontsize=9.5, fontcolor="#7B8794"];
%   rankdir=TB;
%   Claim [label=<<b>1. Claim slot</b><br/><font point-size="9" color="#2F6B58">INSERT ... ON CONFLICT DO NOTHING</font>>, fillcolor="#B7DDD0", color="#6FA890", fontcolor="#1F4E39"];
%   Skip  [label=<<b>Skip</b><br/><font point-size="9" color="#6B2A2A">another worker owns the slot</font>>, fillcolor="#F6C6C6", color="#D98C8C", fontcolor="#6B2A2A"];
%   Run   [label=<<b>2. Mark RUNNING</b><br/><font point-size="9" color="#2F6B58">only from SCHEDULED</font>>, fillcolor="#B7DDD0", color="#6FA890", fontcolor="#1F4E39"];
%   Route [label=<<b>3. Resolve internal route</b><br/><font point-size="9" color="#41719C">contract must be ACTIVE; no buyer key</font>>, fillcolor="#D3E3F3", color="#7CA6CE", fontcolor="#1F4E79"];
%   Build [label=<<b>4. Build provider call</b><br/><font point-size="9" color="#41719C">one user message, temperature 0, timeout 2 × L<SUB>max</SUB></font>>, fillcolor="#D3E3F3", color="#7CA6CE", fontcolor="#1F4E79"];
%   Call  [label=<<b>5. Call seller provider</b><br/><font point-size="9" color="#41719C">through the demo fault hook</font>>, fillcolor="#D3E3F3", color="#7CA6CE", fontcolor="#1F4E79"];
%   Seller [label=<<b>Seller provider</b>>, fillcolor="#E9E1F7", color="#A88FD9", fontcolor="#45296B"];
%   Ok    [label="Status ok?"];
%   Grade [label=<<b>6a. Grade reply</b><br/><font point-size="9" color="#41719C">correct is 1 or 0, or unknown</font>>, fillcolor="#D3E3F3", color="#7CA6CE", fontcolor="#1F4E79"];
%   NoGrade [label=<<b>6b. No grade</b><br/><font point-size="9" color="#41719C">correct is unknown, reason PROVIDER_status</font>>, fillcolor="#D3E3F3", color="#7CA6CE", fontcolor="#1F4E79"];
%   Log   [label=<<b>7. Log request row</b><br/><font point-size="9" color="#2F6B58">is_canary, canary_id, canary_correct, grader_version</font>>, fillcolor="#B7DDD0", color="#6FA890", fontcolor="#1F4E39"];
%   Done  [label=<<b>8. Mark COMPLETED</b><br/><font point-size="9" color="#2F6B58">link request, keep grader version</font>>, fillcolor="#B7DDD0", color="#6FA890", fontcolor="#1F4E39"];
%   Fail  [label=<<b>FAILED_INTERNAL</b><br/><font point-size="9" color="#6B2A2A">error stored, later probes continue</font>>, fillcolor="#F6C6C6", color="#D98C8C", fontcolor="#6B2A2A"];
% 
%   Claim -> Run [label="New row"];
%   Claim -> Skip [label="Conflict", style=dashed];
%   Run -> Route;
%   Route -> Build;
%   Build -> Call;
%   Call -> Seller [dir=both, penwidth=1.6, color="#C0455B"];
%   Call -> Ok;
%   Ok -> Grade [label="Yes"];
%   Ok -> NoGrade [label="No"];
%   Grade -> Log;
%   NoGrade -> Log;
%   Log -> Done;
%   Run -> Fail [style=dashed, label="Exception in steps 2 to 8"];
%   { rank=same; Claim; Skip; }
%   { rank=same; Call; Seller; }
%   { rank=same; Grade; NoGrade; }
% }
% ```
% label=fig:probe-execution
% caption=Execution of one claimed probe slot. Teal steps change database state, blue steps are processing, violet is the external seller, and coral marks the skip and failure paths.
% rendered_images:end
% render_images:begin
\begin{figure}[H]
  \includegraphics[width=\linewidth]{06_probing.tex.figs/06_probing.2.png}
  \caption{Execution of one claimed probe slot. Teal steps change database state, blue steps are processing, violet is the external seller, and coral marks the skip and failure paths.}
  \label{fig:probe-execution}
\end{figure}
% render_images:end

The request that the seller receives is built by \texttt{build\_probe\_call}.
It contains the contract's model identifier, a single user message with the
canary prompt, an output limit of at most 64 tokens (capped by the gateway
setting), temperature $0$, and a timeout of $2 \times L_{\max}$, where
$L_{\max}$ is the contract latency. It carries no field that identifies a
probe. The canary identifier, the grade, and the grader version exist only in
the platform's own \texttt{requests} row, which is written after the response
returns. Table~\ref{tab:visibility} lists what each party can observe.

\begin{table}[htbp]
  \centering
  \small
  \begin{tabularx}{\linewidth}{@{}p{0.32\linewidth}XX@{}}
    \toprule
    Item & Seller provider & Platform database \\
    \midrule
    Prompt text & visible & stored in \texttt{requests} \\
    Model, output limit, temperature & visible & stored \\
    \texttt{is\_canary}, \texttt{canary\_id} & not sent & stored \\
    Expected answer and grade & not sent & stored \\
    Probe slot and scheduled time & not sent & stored in
      \texttt{canary\_probes} \\
    Gateway route & pinned provider slug & pinned and actual provider \\
    \bottomrule
  \end{tabularx}
  \caption{Visibility of probe metadata. The seller receives the same message
  fields as for buyer traffic, apart from the prompt distribution.}
  \label{tab:visibility}
\end{table}

The probe uses the gateway's normal serving path through an internal route
resolver. The function \texttt{resolve\_internal\_route} loads the contract and
seller by identifier, requires the contract to be \texttt{ACTIVE}, and returns
the same route object as the buyer-facing resolver, but without buyer
authentication, because the prober acts as the platform. The call passes
through the demo fault hook, so a fault injected for a seller affects probes as
well as buyer traffic. The grade, or a null grade for a failed provider call,
is stored with the request together with the grader version, and the probe row
is then completed with a link to the request.

% ==============================================================================
\subsection{Concurrency and failure semantics}
\label{sec:failure}

Several workers can run against one database. A slot is claimed with
\texttt{INSERT ... ON CONFLICT (contract\_id, probe\_slot) DO NOTHING}, and only
the worker that inserted the row proceeds, which gives at-most-once execution
per slot. State transitions are guarded (\texttt{SCHEDULED} to
\texttt{RUNNING}, \texttt{RUNNING} to \texttt{COMPLETED}), so a second
attempt on the same row fails visibly. The worker polls once per second by
default and stops after the current pass on \texttt{SIGINT} or \texttt{SIGTERM}.
A failed database-level pass is logged and retried.

Table~\ref{tab:failures} lists how failures are recorded and whether the seller
is charged. The principle is P5: an outcome counts against the seller only if
the provider produced it.

\begin{table}[htbp]
  \centering
  \small
  \begin{tabularx}{\linewidth}{@{}p{0.3\linewidth}Xp{0.17\linewidth}@{}}
    \toprule
    Event & Recorded as & Charged to seller \\
    \midrule
    Provider timeout or upstream error & request with status \texttt{timeout}
      or \texttt{upstream\_error}, no grade & yes \\
    Wrong or malformed answer & request with grade 0 & yes \\
    Gateway or configuration error & request with status
      \texttt{gateway\_error} & no (excluded) \\
    Grader cannot grade a successful reply & grade $\bot$ & no (checker
      error) \\
    Exception inside a probe & probe \texttt{FAILED\_INTERNAL} with the error
      text; later probes continue & no \\
    Scheduling failure for one contract & logged and returned in the cycle
      result; other contracts continue & no \\
    Slot already claimed & skipped & no \\
    \bottomrule
  \end{tabularx}
  \caption{Failure handling in the prober.}
  \label{tab:failures}
\end{table}

\paragraph{Operational notes.} The prober code does not reclaim a probe that
remains \texttt{RUNNING} after a worker crash, so the unique slot key then
prevents a retry of that slot. The effect is a missing observation, which
lowers the evidence count and can only move a verdict toward
\texttt{INSUFFICIENT}. Probes within a pass run sequentially. Section~\ref{sec:limitations}
discusses these points.
